\section{Baseline Explanation: ForeAgent}
\label{sec:baseline}

\subsection{ML-Agent Task and Execution Feedback}
Given a natural-language instruction $I$ and dataset $D$, an autonomous ML agent seeks the optimal code solution $C^*$ in code space $\mathcal C$ by maximizing task metric $M$ \citep[Section~2.1]{zheng2026predict}:
\begin{equation}
 C^*=\arg\max_{C\in\mathcal C} M(I,C,D).
 \label{eq:agent-objective}
\end{equation}
Here, $C$ is a candidate program, $M(I,C,D)$ measures its task performance, and $C^*$ is an optimal program. Under a fixed compute budget, the agent can execute and evaluate only a subset of $\mathcal C$.

\newpage
\subsection{Data-centric Solution Preference and the Corpus}
\label{sec:preference}
\citet{zheng2026predict} define Data-centric Solution Preference as a pairwise selection task. Given a task description, a data report, two candidate code solutions, and a system prompt, an LLM judge reasons about which candidate is likely to perform better before either candidate is executed. Its input and output are
\begin{equation}
 \begin{aligned}
 X&=\left(I,D_{\mathrm{rep}},\{C_0,C_1\},P\right),\\
 \mathcal{Y}&=\left\{\left(\cot,\hat y,c\right)\mid \cot,\hat y\in\{0,1\},\;c\in[0,1.0]\right\}.
 \end{aligned}
 \label{eq:predictor}
\end{equation}
Here, $I$ is the task description; $D_{\mathrm{rep}}$ is the authors' Verified Data Analysis Report, supplied as context about the task data; $C_0$ and $C_1$ are the candidate code solutions; and $P$ is the system prompt. The output contains $\cot$, the judge's reasoning; $\hat y\in\{0,1\}$, the predicted winner (0 for $C_0$ and 1 for $C_1$); and $c\in[0,1]$, the judge's reported confidence. The judge's prediction is distinct from the ground-truth preference, which is assigned using the candidates' recorded execution scores. ForeAgent uses the LLM judge without training its parameters.

To evaluate this predictive task, the authors collected solutions generated by AIDE and AutoMind on 26 MLE-bench tasks adapted from Kaggle competitions \citep{zheng2026predict,chan2024mle}. They curated 895 solutions from 1,329 valid candidates and constructed 18,438 within-task pairs, with preference labels derived from execution results. On this corpus, DeepSeek-V3.2-Thinking and GPT-5.1 achieved 61.5\% and 58.8\% pairwise accuracy, respectively. These results provide evidence that an LLM's pre-execution judgments contain a useful, though imperfect, signal about candidate performance. The released execution records also provide a potential resource for future reward-model training.

\subsection{System Architecture and Design Choices}
ForeAgent applies these judgments within AIDE's code-search framework. AIDE stores candidate programs as nodes in a solution tree and uses execution feedback to decide whether to draft a new program, debug a failed one, or improve a valid one. Its journal retains earlier programs, plans, errors, and measured results, allowing later code-generation prompts to draw on the search history \citep{jiang2025aide}. At revision step $t$, the general execution-feedback update is \citep[Section~2.2]{zheng2026predict}
\begin{equation}
 C_{t+1}\leftarrow\operatorname{Agent}(I,C_t,
 \underbrace{\operatorname{Execute}(C_t,D)}_{R_t})
 \label{eq:execution-feedback}
\end{equation}
Here, $C_t$ is the selected code at step $t$, $R_t$ is its execution feedback (such as errors or validation results), and $C_{t+1}$ is the next proposed code. ForeAgent retains AIDE's framework and redesigns its improvement stage as a Predict-then-Verify loop. Drafting and debugging still lead to direct execution; improvement first generates and compares several possible revisions before committing execution resources. Figure~2(e) of \citet{zheng2026predict} illustrates this branch of the workflow.

The loop has three phases. First, candidate generation: when the controller chooses to improve a valid parent program, the coding model receives the task description, the parent code, and a journal summary of earlier attempts. It produces an improvement plan and corresponding Python code for up to $m=10$ revisions. The prompts ask for a specific change and encourage distinct improvement directions across the batch. Second, confidence-gated selection: the comparison model judges pairs of proposed programs using the input defined in Section~\ref{sec:preference}. Third, verification execution: selected code is run, and its errors and measured validation result are written back to the journal for subsequent search steps. In the authors' experimental setup, DeepSeek-V3.2 generates code and DeepSeek-V3.2-Thinking serves as the comparison model \citep[Sections~6.2--6.3]{zheng2026predict}.

\subsection{Inference-Time Selection and Execution Rules}
\label{sec:selection-rules}
ForeAgent's comparison model predicts a winner and confidence for each code pair. For a comparison between candidates $C_i$ and $C_j$, write this confidence as $c_{ij}$. The released configuration sets the confidence threshold to $\tau=0.7$. A judgment with $c_{ij}\geq\tau$ advances its predicted winner and eliminates the loser. When confidence falls below the threshold, the candidates remain unresolved and receive further comparison. The author implementation organizes these decisions as a tournament (Algorithm~\ref{alg:tournament}).

\begin{algorithm}[H]
\caption{Confidence-gated tournament (released implementation)}\label{alg:tournament}
\small
\begin{algorithmic}[1]
\Require Improvement candidates $A$, confidence threshold $\tau=0.7$
\State Initialize each candidate's win count and elimination status
\While{multiple candidates remain, subject to a round limit}
\State Pair adjacent candidates; advance an unpaired candidate
\State Compare each pair with the LLM judge
\For{each pair}
\If{confidence $\geq\tau$}
\State Advance the predicted winner and eliminate the loser
\Else
\State Place both candidates in a pending group
\EndIf
\EndFor
\For{each pending candidate}
\State Compare it with a randomly chosen advanced candidate
\If{it loses with confidence $\geq\tau$}
\State Eliminate it
\Else
\State Advance it
\EndIf
\EndFor
\State Continue with the advanced candidates
\EndWhile
\State Rank candidates by survival, win count, and round reached
\State \Return the ranking
\end{algorithmic}
\end{algorithm}

The paper summarizes verification as selecting the top $k$ candidates, with $k=1$. The released implementation \citep{foreagentcode} adds an execution rule: each eligible rank-zero candidate is executed with probability $p_{\mathrm{exec}}=0.2$; if every eligible candidate is skipped, the agent executes one rank-zero candidate as a fallback. Candidates can share rank zero, so $k=1$ can make more than one candidate eligible. Execution trains and evaluates the selected candidate program, after which the agent records its feedback and continues the search. The final program is chosen from executed, valid solutions using the measured validation metric.

ForeAgent does not train its LLM judge as part of this procedure. Candidate programs may train task-specific ML models when executed, using losses specified by their generated code; ForeAgent itself is governed by its pairwise predictions, confidence threshold, tournament ranking, and execution rule. These are therefore the baseline's relevant inference-time decision and scoring rules, rather than a single training loss for the agent.

\subsection{Reported Evaluation}
Zheng et al.\ evaluated ForeAgent against AIDE on five MLE-bench tasks, including two designated unseen tasks. Each task had a 12-hour limit and was run three times. Their primary outcome, Beat Ratio, measures the proportion of human leaderboard contestants whose scores the agent surpasses. The authors report an average Beat Ratio of 0.739 for ForeAgent versus 0.695 for AIDE, approximately a 6\% relative improvement. They also report reaching peak performance with about one-sixth of the execution time and exploring 3.2 times as many nodes on average \citep[Sections~6.3--6.4]{zheng2026predict}. These are the original paper's five-task results; the two-task, single-run reproduction is presented separately in Section~\ref{sec:reproduction}.
